        Here we present an architectural comparison between our proposed Multi-Output Temporal SR with Learned Up-scaling (MO-TSR) model and state-of-the-art SR methods. Methods are organized into four categories: CNN-based, GAN-based, Transformer-based, and satellite multi-frame specific methods. The comparison was made across six dimensions; namely, feature extraction, attention mechanism, up-sampling strategy, loss functions, temporal and multi-frame handling, and multi-resolution output characteristics.

        In terms of feature extraction, MO-TSR uses a dual-resolution encoder with separate branches for 10\,m and 20\,m Sentinel-2 bands, each composed of a $7 \times 7$ convolutional stem followed by four residual blocks with 64 base channels and CBAM (Convolutional Block Attention Module) applied at alternating blocks. This dual-resolution encoder is not present in any of the SOTA models, representing an innovation and granting SR algorithms the capacity to inherently handle multi-resolution bands. HighRes-net \citep{deudon2020highresnetrecursivefusionmultiframe} uses a shared temporal fusion encoder that allows it to merge same-resolution frames over time; however, MO-TSR's dual encoder allows downstream applications not only to process multi-frame inputs but also to build representations from different spectral-spatial resolution bands. This feature advances the SR field for satellite imagery considering Sentinel-2 has varying spatial resolutions across its bands and it is commonly used as the LR source as it offers radiometrically consistent medium-resolution open-source imagery at a global scale. 
        
        Other unique features of the proposed MO-TSR are the temporal cross-attention mechanism and the gated cross-resolution fusion. The temporal cross-attention module is the core of the temporal fusion block and creates a dual path for global and local temporal information sharing. The global path applies scaled dot-product attention ($QK^T/\sqrt{d}$) over spatially-pooled frame representations, producing a frame-level attention matrix ($T \times T$, where $T = 3$) that captures inter-frame relevance. The local path concatenates all temporal frames along the channel dimension and applies a $1 \times 1$ convolution to produce a spatially-varying local mixture. Per-frame outputs are then computed as a gated interpolation between each frame's own features and the local mixture, with an additive global residual:
        %
        \begin{equation}
            \mathbf{o}_t = (1 - \mathbf{g}_t) \odot \mathbf{f}_t + \mathbf{g}_t \odot \mathbf{m}_{\text{local}} + 0.1 \cdot \mathbf{v}_t^{\text{global}},
        \end{equation}
        %
        where $\mathbf{g}_t = \sigma(W[\mathbf{f}_t \| \mathbf{m}_{\text{local}}])$ is a learned sigmoid gate. The gated cross-resolution fusion mechanism then creates a per-channel, spatially-varying, bidirectional information exchange between 10\,m and 20\,m feature maps. Specifically, the 20\,m features are bilinearly upsampled to 10\,m spatial resolution (and vice versa), transformed through learned $1 \times 1$ convolutions, and combined via a gated residual addition:
        %
        \begin{equation}
            \mathbf{f}_{10}^{\prime} = \mathbf{f}_{10} + \sigma\!\left(W\!\left[\mathbf{f}_{10} \| T_{20 \to 10}(\mathbf{f}_{20})\right]\right) \odot T_{20 \to 10}(\mathbf{f}_{20}),
        \end{equation}
        %
        and symmetrically for the 20\,m branch. The 20\,m bands provide additional spectral information for the 10\,m bands, while the 10\,m bands help inform the 20\,m bands with additional textural information. MO-TSR uses a spectral consistency channel attention module to preserve spectral coherence that can be distorted from multi-spectral band information sharing. This module applies adaptive average pooling followed by two $1 \times 1$ convolutions with a GELU activation and a sigmoid gate, recalibrating channel responses to maintain inter-band relationships. MO-TSR shares a Squeeze-and-Excitation (SE) channel attention strategy with other SOTA models such as RCAN and RAMS, which is useful to let the SR model focus on detailed information in the image. This can be associated to the observed capacity of both the direct prediction and residual learning versions of the MO-TSR to learn high-frequency features.

        One major limitation to train satellite SR algorithms with real datasets is the up-scaling factor that exists between LR and HR images. For instance, the up-scaling factor between Sentinel-2 and PlanetScope is $3.33\times$ for both the 10\,m bands (10\,m $\to$ 3\,m GSD) and the 20\,m bands (20\,m $\to$ 6\,m GSD). This non-integer up-scaling factor, which also exists with other HR image providers (e.g., WorldView), limits the applicability of the PixelShuffle up-scaling mechanism that is used by almost every DL-based SR algorithm. In MO-TSR we use a progressive hybrid approach that combines PixelShuffle ($2\times$) with Adaptive Interpolation. First, PixelShuffle doubles the spatial resolution (e.g., $78 \to 156$ for 10\,m bands, $39 \to 78$ for 20\,m bands), followed by a $3 \times 3$ output convolution. Then, the Adaptive Interpolation module---a content-aware deformable $5 \times 5$ kernel with learned offsets bounded by $\tanh(\cdot) \times 2.0$ pixels---resizes the intermediate features to the exact target size (e.g., $156 \to 260$ for 10\,m, $78 \to 130$ for 20\,m), achieving the remaining $1.667\times$ factor. This allows the SR model to handle non-integer up-scaling factors and prevents over-smoothing of high-frequency regions that is common in traditional non-DL SR methods \citep{Zhang2022223}. While the former advantage is essentially irrelevant for non-satellite applications, for satellite imagery SR, preserving geometric fidelity and pixel information is essential for many downstream applications such as spectral indices calculation and change detection. Other satellite SR applications still choose to use larger integer up-scaling factors and then center-crop to the desired size. This strategy introduces edge artifacts at crop boundaries whose relevance depends on the end application. Some limitations of this approach include the larger computational cost (compared to only using PixelShuffle) of the adaptive interpolation, which cannot be parallelized due to the per-pixel offset computation, and the memory overhead when used with large input patches.

        Finally, one major incorporation used in the development of MO-TSR is the compound domain-specific loss. Most SOTA SR methods only use single or few loss components---for instance, EDSR, RCAN, HAT, SRFormer, DAT, and MambaIR all employ only L1 loss, while HighRes-net and RAMS similarly rely on L1 alone. Only a few SR models such as ESRGAN \citep{wang2018esrganenhancedsuperresolutiongenerative}, and its variants, and SwinIR \citep{liang2021swinirimagerestorationusing} use a compound loss term combining L1, VGG perceptual, and adversarial losses; however, none incorporates SAM or temporal loss terms in their model training routine. MO-TSR employs a four-component compound loss: (1) L1 pixel reconstruction loss ($\alpha = 1.0$), (2) Spectral Angle Mapper (SAM) loss ($\alpha = 0.137$) computed as the mean arccos of the normalized cosine similarity across spectral bands, (3) local gradient loss ($\alpha = 0.01$) using Sobel-filtered edge maps to preserve spatial detail, and (4) temporal consistency loss ($\alpha = 0.5$) that penalizes differences in temporal change patterns between predicted and target sequences. Incorporating multiple loss terms beyond L1 has been studied and proven to be beneficial, considering that there is a non-unique LR-HR mapping path that is overlooked by the L1 loss term \citep{he2023revisitingl1losssuperresolution}. The SAM and temporal consistency terms are particularly relevant for remote sensing: the SAM loss explicitly preserves inter-band spectral relationships critical for applications such as vegetation index computation, while the temporal consistency loss encourages coherent predictions across the $T = 3$ frame sequence, a constraint absent from all single-image SOTA methods.
